\subsection{A smooth fixed divisor and a singular quotient: Problem 14.72}
\label{cand:14.72}

Over \(\C\), let \(X\subset\mathbb A^4\) have coordinates
\((x,y,u,v)\) and equations
\[
 y^2=x^3-x,\qquad yu=xv,\qquad (x^2-1)u=yv.
\]
The automorphism
\(\sigma(x,y,u,v)=(-x,iy,iu,v)\) has order four.
The curve \(C:y^2=x^3-x\) is smooth and irreducible.
The opens \(x\ne0\) and \(x^2-1\ne0\) cover \(X\).
On the first, eliminate \(v=yu/x\); on the second,
eliminate \(u=yv/(x^2-1)\). They identify \(X\)
locally with the product of an open subset of \(C\)
and an affine line, proving it is a smooth irreducible
affine surface without hidden components.

The full fixed scheme of \(\langle\sigma\rangle\) is
\(x=y=u=0\), with \(v\) free. It is a smooth affine
line. On the second chart its ideal is generated by
\(y\), because \(x=y^2/(x^2-1)\) and
\(u=yv/(x^2-1)\). It is therefore a smooth effective
Cartier divisor, of the codimension one required in
the question.

Nevertheless the quotient is singular. Put
\[
 a=x^2,\qquad d=xu^2-v^2.
\]
These and \(v\) are invariant, and the defining equations
give \(xd=u^2\) and \((a-1)d=v^2\). We claim
\[
 \C[X]^{\langle\sigma\rangle}
       =\C[a,v,d]/(v^2-(a-1)d).
\]
After inverting \(a\), a basis of the coordinate ring
consists of \(x^s y^b u^c\), \(s\in\Z\), \(b=0,1\),
\(c\geq0\), with weight \(i^{2s+b+c}\).
The invariant monomials have respectively the forms
\[
 x^s u^{2h}=a^{(s+h)/2}d^h,\qquad
 x^s y u^{2h+1}=a^{(s+h+1)/2}vd^h.
\]
Invariance makes the displayed exponents integral.
The monomials \(a^r d^h,a^r v d^h\) have distinct
images in that Laurent basis, so there are no
further relations. After inverting \(a-1\), the other
chart has invariant ring
\(\C[a,(a-1)^{-1},v]\), since the invariant ring of
the curve is \(\C[x^2]\). Here \(d=v^2/(a-1)\).
The two invariant principal opens cover, so the
identifications prove the global assertion. One can
also clear powers of \(a\) and \(a-1\) and use their
comaximality to see that no invariant is omitted.

At \((a,v,d)=(1,0,0)\) the local equation is
\(v^2=(a-1)d\). The local dimension is two and the
cotangent dimension three, so the quotient is not
regular. This point comes from the orbit of
\((1,0,0,0)\), whose stabilizer is the proper subgroup
\(\langle\sigma^2\rangle\). The full fixed divisor
does not control this additional stabilizer.
The construction uses cyclic order four; it makes
no assertion about a prime-order variant or an action
free outside the fixed divisor. Source:
\artifact{research/14.72-proof.md}.

